\documentclass[12pt]{article}
\usepackage[utf8]{inputenc}
\usepackage{listings}
\usepackage{xcolor}
\usepackage{geometry}
\usepackage{parskip}

\geometry{margin=2.5cm}

\definecolor{codebg}{RGB}{245,245,245}
\definecolor{keyword}{RGB}{0,112,0}
\definecolor{comment}{RGB}{128,128,128}
\definecolor{string}{RGB}{180,0,0}
\definecolor{outputbg}{RGB}{240,248,255}

\lstdefinestyle{pythonstyle}{
    backgroundcolor=\color{codebg},
    basicstyle=\ttfamily\small,
    keywordstyle=\color{keyword}\bfseries,
    stringstyle=\color{string},
    commentstyle=\color{comment}\itshape,
    language=Python,
    frame=single,
    rulecolor=\color{gray!40},
    breaklines=true,
    showstringspaces=false,
    tabsize=4,
    captionpos=b,
}

\lstdefinestyle{outputstyle}{
    backgroundcolor=\color{outputbg},
    basicstyle=\ttfamily\scriptsize,
    frame=single,
    rulecolor=\color{blue!20},
    breaklines=true,
}

\title{\textbf{Wind Turbine Blade Fault Classification\\Using a Simple CNN}\\[0.5em]
\large Step 9 --- Train the Model}
\author{Amreen Batool}
\date{}

\begin{document}
\maketitle

\section*{Step 9 --- Train the Model}

\begin{lstlisting}[style=pythonstyle]
early_stop = EarlyStopping(
    monitor='val_loss',
    patience=3,
    restore_best_weights=True
)

history = model.fit(
    train_data,
    validation_data=val_data,
    epochs=10,
    callbacks=[early_stop]
)
\end{lstlisting}

\subsection*{Output}

\begin{lstlisting}[style=outputstyle]
Epoch 1/10  263/263 - accuracy: 0.5573 - loss: 0.6943 - val_accuracy: 0.6233 - val_loss: 0.6348
Epoch 2/10  263/263 - accuracy: 0.6201 - loss: 0.6432 - val_accuracy: 0.6350 - val_loss: 0.6328
Epoch 3/10  263/263 - accuracy: 0.6309 - loss: 0.6283 - val_accuracy: 0.5950 - val_loss: 0.6410
Epoch 4/10  263/263 - accuracy: 0.6409 - loss: 0.6186 - val_accuracy: 0.6533 - val_loss: 0.6285
Epoch 5/10  263/263 - accuracy: 0.6613 - loss: 0.6056 - val_accuracy: 0.6750 - val_loss: 0.6075
Epoch 6/10  263/263 - accuracy: 0.6724 - loss: 0.6027 - val_accuracy: 0.6983 - val_loss: 0.5723
Epoch 7/10  263/263 - accuracy: 0.6557 - loss: 0.5974 - val_accuracy: 0.6683 - val_loss: 0.6180
Epoch 8/10  263/263 - accuracy: 0.6705 - loss: 0.5916 - val_accuracy: 0.6900 - val_loss: 0.5746
Epoch 9/10  263/263 - accuracy: 0.6785 - loss: 0.5757 - val_accuracy: 0.6717 - val_loss: 0.5590
Epoch 10/10 263/263 - accuracy: 0.7038 - loss: 0.5525 - val_accuracy: 0.7167 - val_loss: 0.5334
\end{lstlisting}

\subsection*{Explanation}

\textbf{\texttt{EarlyStopping}:}
\begin{itemize}
    \item Monitors the \texttt{val\_loss} (validation loss) after each epoch.
    \item If validation loss does not improve for 3 consecutive epochs (\texttt{patience=3}), training stops automatically.
    \item \texttt{restore\_best\_weights=True} ensures the model keeps the weights from the best epoch.
\end{itemize}

\textbf{\texttt{model.fit(...)}:}
\begin{itemize}
    \item Trains the model on the training data.
    \item \texttt{epochs=10} --- the model can train for a maximum of 10 rounds, but early stopping may halt it sooner.
    \item \texttt{validation\_data=val\_data} --- performance is checked on the validation set after every epoch.
\end{itemize}

\subsection*{Why This Is Good}

Early stopping keeps the model from overtraining on the training data, which helps maintain better performance on unseen test images.

\end{document}
